\section{Diseño del mapeo visual}
El diseño de la representación tridimensional tiene como objetivo transformar las estructuras 
abstractas y los tensores multidimensionales de la red neuronal en primitivas geométricas e 
interactivas dentro de Blender. Para lograr una interpretación intuitiva del modelo, el mapeo 
visual se articula en tres dominios fundamentales: la \textbf{geometría de las capas}, la 
\textbf{topología de la red} y la \textbf{reconstrucción de la salida}.

\subsection{Mapeo de capas y tensores de activación}
Cada capa convolucional o de agrupamiento (\textit{Max-Pooling}) se representa como un 
paralelepípedo contenedor (\textit{bounding volume}), cuyas dimensiones físicas de ancho, 
alto y profundidad en el espacio 3D son proporcionales a la resolución espacial y número de 
canales del tensor correspondiente.
\medskip

Dentro de estos contenedores, la información volumétrica proveniente de los archivos 
\textit{OpenVDB} se renderiza mediante un esquema de rebanado dinámico (\textit{planar slicing}). 
Este mecanismo permite al usuario interpolar a través de la profundidad del tensor para inspeccionar 
mapas de activación específicos. En ausencia de datos volumétricos o con fines de abstracción 
arquitectónica, el contenedor puede representarse como un cuerpo sólido simplificado para reducir 
la carga de renderizado. Asimismo, los metadatos de la capa (identificador, tipo de operación, 
dimensiones del tensor y tamaño del \textit{kernel}) se proyectan en la escena mediante etiquetas
de texto 3D vinculadas a la posición de la capa.

\subsection{Mapeo de la topología y flujo de datos}
La topología de la red se formaliza mediante un grafo dirigido en el espacio tridimensional. 
Las conexiones entre capas consecutivas se representan mediante cuatro aristas geométricas que unen 
las esquinas correspondientes de sus volúmenes delimitadores; específicamente, conectan la cara 
posterior del paralelepípedo de la capa de entrada con la cara frontal del paralelepípedo de la capa 
de salida. Esta representación en cuadrilátero proyectado define un volumen de transición claro que 
delimita visualmente el flujo de datos entre capas.
\medskip

Para representar el flujo de procesamiento, el sistema modela el comportamiento del kernel 
como una primitiva geométrica secundaria que se desplaza o proyecta 
iterativamente sobre las rebanadas del tensor, evidenciando las relaciones de dependencia espacial y 
las operaciones de submuestreo entre capas conectadas.

\subsection{Mapeo de la salida y predicciones}
La representación de la salida del modelo integra la imagen de entrada en el plano focal de la escena 
y superpone las predicciones geométricas filtradas. Los valores de confianza (\textit{confidence}) y 
probabilidad de clase (\textit{class probability}) extraídos por el modelo no poseen una primitiva 
geométrica directa; en su lugar, para filtrar y 
discriminar las cajas que superan los umbrales configurados, ocultando de la visualización las 
detecciones no significativas.
\medskip

Las cajas delimitadoras (\textit{bounding boxes}) resultantes del filtrado se instancian sobre la 
imagen como mallas alámbricas a las que se les aplican las transformaciones geométricas de posición y 
escala requeridas según sus parámetros $(x, y, w, h)$. Finalmente, la categoría o clase detectada 
se visualiza mediante dos elementos explícitos: el color asignado a la estructura de la caja y 
una malla de texto 3D que despliega la etiqueta correspondiente sobre el volumen delimitador.

\subsection{Resumen de primitivas visuales}
En la Cuadro \ref{tab:mapeo_visual} se sintetiza la correspondencia entre los elementos de la red 
neuronal y sus primitivas de representación gráfica en el entorno 3D.

\begin{table}[htbp]
    \centering
    \small
    \begin{tabular}{lll}
        \hline
        \textbf{Elemento de la Red} & \textbf{Atributo / Datos} & \textbf{Primitiva / Representación en Blender} \\ \hline
        Capas & Dimensiones del tensor & Paralelepípedo 3D (Escala proporcional) \\
        Tensores de activación & Salidas volumétricas & Malla OpenVDB con \textit{slicing} de densidad \\
        Metadatos de capa & Nombre, Kernel, Stride & Etiquetas tipográficas \\
        Conexiones inter-capa & Topología de red & 4 aristas entre esquinas de caras opuestas \\
        Operación del Kernel & Paso / Desplazamiento & Primitiva iterativa de barrido \\
        Cajas de predicción & $(x, y, w, h)$ & Malla alámbrica superpuesta (\textit{Bounding Box}) \\
        Confianza y Probabilidad & \textit{Confidence, Class Prob.} & Criterio de filtrado (Sin primitiva directa) \\
        Categoría / Clase & ID de Clase / Etiqueta & Mapeo de color de caja y malla de texto 3D \\ \hline
    \end{tabular}
    \caption{Matriz de mapeo entre los elementos de la red y las primitivas visuales en Blender.}
    \label{tab:mapeo_visual}
\end{table}